\section{Introduction}\label{sec:introduction}

% Chords form an integral part of music. Part of how musicians understand many styles of music is through harmonic structure~\cite{tonalharmony}. 
Chord annotations allow music to be easily shared, performed, improvised and analysed. Creating high-quality chord annotations requires a trained music and can be highly subjective~\cite{annotatorsubjectivity}. Thus, ACR models are a promising avenue for creating high-quality chord annotations at scale, allowing for better understanding, creation and learning of music. To this end, ACR has been tackled by deep learning methods since the influential work of~\cite{RethinkingChordRecognition}. They train a convolutional neural network (CNN) to classify pitch spectra into chord labels. Progress was made by combining CNNs with recurrent neural networks (RNNs)~\cite{firstrnn, ACRCNNRNN1,ACRLargeVocab1,StructuredTraining} with a CNN performing feature extraction from a spectrogram and an RNN sharing information across frames. Such models are referred to as CRNNs. More recently, transformers have been applied in place of the RNN~\cite{MelodyTranscriptionViaGenerativePreTraining, HarmonyTransformer, AttendToChords,CurriculumLearning,BTC,discordharmony,chordformer}. Despite increasingly complex models being employed for ACR, performance improvements have stagnated, and have barely improved in 10 years since the first CRNN trained by~\cite{firstrnn}. The `glass ceiling' discussed by~\cite{FourTimelyInsights} is still present 10 years later.


The original work in~\cite{RethinkingChordRecognition} classified pitch spectra into 25 chords labels — all 12 keys of minor and major plus a no chord symbol. Most work since then has focused on a much larger vocabulary. The most common is the 170 chord vocabulary introduced by~\cite{StructuredTraining}, which includes 14 chord qualities, plus a no chord and unknown chord symbols. Such a large chord vocabulary creates a long tail of rare chord labels. For example, major chords are very common, while minmaj7 chords are extremely rare our dataset. Detailed chord annotations can also be highly subjective~\cite{annotatorsubjectivity}. For example, annotators may disagree on the presence of a chord extension. Such subjectivity may make it harder for a model to learn a consistent mapping from audio to chord labels, and brings into question the upper bound of performance for ACR models.

In this work, we propose that among the main reasons for the lack of performance improvement in ACR is due to the limitation imposed by the size of the dataset. The most common dataset used for ACR — and the dataset used in this work and other works over the last decade — contains around 1200 annotated pop songs. By contrast, Jukebox~\cite{Jukebox}, a generative music model published over 5 years ago, was trained on 1.2 million songs. It is infeasible to obtain an annotated dataset of any comparable size for ACR. Instead, we must find ways of leveraging self-supervised methods, models from other musical domains or generative models to learn from a larger source of data, without more chord labels. Here, we propose three methods: (1) extracting \emph{generative features} from MusicGen~\cite{MusicGen} as input to a CRNN (2) using a chord-conditioned generative model to create synthetic data to train on and (3) using a beat estimation model to inform the model of when chord changes are likely to occur. We perform our experiments by reimplementing the CRNN from~\cite{StructuredTraining} and ablating on various common performance improvements. This includes output smoothing, structured decomposition of the chord vocabulary, and pitch augmentation. We then apply the three methods described above and find that none meaningfully improve performance. However, training on synthetic data does improve performance marginally, and changes the types of errors made by the model. We also find that with oracle beat estimation, performance improves significantly. All this suggests that better synthetic datasets and beat estimation models could be a promising avenue for improving ACR performance and breaking through the current `glass ceiling'. All code and a more detailed manuscript containing all experimental results are available on GitHub.\footnote{redacted for anonymity} Data can be made available upon request.\footnote{redacted for anonymity}

% The remainder of paper is structured as follows. \textbf{Section~\ref{sec:related-work}} discusses existing literature and the potential for foundation models to help. \textbf{Section~\ref{sec:method}} describes the datasets, evaluation metrics, models and training procedure used, as well as a large series of ablations on important results from previous works. \textbf{Section~\ref{sec:experiments}} describes the experiments and results when using generative features, synthetic data and beat-estimation models. \textbf{Section~\ref{sec:conclusion}} concludes the paper and provides suggestions for future work. \textbf{Section~\ref{sec:extras}} contains acknowledgments, limitations, an AI usage statement and, as is increasingly important for all machine learning research, an ethics statement. All code is available on GitHub.\footnote{\url{https://github.com/PierreRL/AutomaticChordRecognition}} Data can be made available upon request.\footnote{lardet[dot]pierre[at]gmail.com}